\documentclass[letterpaper,11pt,reqno]{amsart} 
\usepackage[portrait,margin=3cm]{geometry}
% Existing personal style, inlined for the standalone editor.
\usepackage{systeme}
\usepackage{mathrsfs,xfrac} 
\usepackage[colorlinks=true,linkcolor=blue,citecolor=blue,urlcolor=blue]{hyperref} 
\usepackage{amsmath,amssymb,amsthm,amsfonts,amsbsy,latexsym,dsfont,color} 
\usepackage[numeric,initials,nobysame]{amsrefs} 
\usepackage[foot]{amsaddr}
%\usepackage{enumitem}
%\usepackage{upref,halloweenmath}
%\renewcommand{\skull}{\vcenter{\hbox{\scalebox{.85}{$\bigskull$}}}}

\usepackage[utf8]{inputenc}
\usepackage[english]{babel}
\usepackage{comment}

\usepackage[textsize=tiny]{todonotes}
\usepackage{regexpatch}
%\tracingxpatches%for debugging
\makeatletter
\xpatchcmd{\@todo}{\setkeys{todonotes}{#1}}{\setkeys{todonotes}{inline,#1}}{}{}
\makeatother

\usepackage{enumerate}
\newenvironment{enumeratei}{\begin{enumerate}[\upshape (i)]}{\end{enumerate}}
\newenvironment{enumeratea}{\begin{enumerate}[\upshape (a)]}{\end{enumerate}}
\newenvironment{enumeraten}{\begin{enumerate}[\upshape 1.]}{\end{enumerate}}
\newenvironment{enumerateA}{\begin{enumerate}[\upshape (A)]}{\end{enumerate}}

%\setlength{\evensidemargin}{.25in}
%\setlength{\textwidth}{6in}
%\setlength{\topmargin}{-0.4in}
%\setlength{\textheight}{8.5in}
%\textwidth=6in
%\oddsidemargin=0.25in
%\evensidemargin=0.25in
%\topmargin=-0.1in
%\footskip=0.8in
%\parindent=0.0cm
%\parskip=0.3cm
%\textheight=8.00in
%\setcounter{tocdepth}{3}
%\setcounter{secnumdepth}{2}
\sloppy

%%%%%%%%%%%%%%%%%%%%%%%%%%%%%%%%%%%%%%%%%%%    
%% Local environments 
\newtheorem{thm}{Theorem}[section]
\newtheorem{lem}[thm]{Lemma}
\newtheorem{cor}[thm]{Corollary}
\newtheorem{prop}[thm]{Proposition}
\newtheorem{defn}[thm]{Definition}
\newtheorem{rem}[thm]{Remark}
\newtheorem{hyp}[thm]{Hypothesis}
\newtheorem{ass}[thm]{Assumption}
\newtheorem{exc}[section]{Exercise}
\newtheorem{ex}[thm]{Example}
\newtheorem{conj}[thm]{Conjecture}   
%% End Local environments
%%%%%%%%%%%%%%%%%%%%%%%%%%%%%%%%%%%%%%%%%%%
%% Local macros
\renewcommand{\le}{\leqslant} 
\renewcommand{\ge}{\geqslant} 
\renewcommand{\leq}{\leqslant} 
\renewcommand{\geq}{\geqslant} 
\newcommand{\eset}{\varnothing}
\newcommand{\ra}{\rangle}
\newcommand{\la}{\langle} 
\newcommand{\wt}{\widetilde}
\newcommand{\pms}{\{-1,1\}}
\newcommand{\ind}{\mathds{1}}
\newcommand{\eps}{\varepsilon}
\newcommand{\To}{\longrightarrow}
\newcommand{\norm}[1]{\left\Vert#1\right\Vert}
\newcommand{\abs}[1]{\left\vert#1\right\vert}
\newcommand{\set}[1]{\left\{#1\right\}}
\newcommand{\goesto}{\longrightarrow}
\newcommand{\Real}{\mathds{R}}
\newcommand{\Complex}{\mathds{C}}
\newcommand{\ie}{\emph{i.e.,}}
\newcommand{\as}{\emph{a.e.}}
\newcommand{\eg}{\emph{e.g.,}}
\renewcommand{\iff}{\Leftrightarrow} 
\newcommand{\equald}{\stackrel{\mathrm{d}}{=}}
\newcommand{\probc}{\stackrel{\mathrm{P}}{\longrightarrow}}
\newcommand{\weakc}{\stackrel{\mathrm{w}}{\longrightarrow}}
\newcommand{\fpar}[2]{\frac{\partial #1}{\partial #2}}
\newcommand{\spar}[2]{\frac{\partial^2 #1}{\partial #2^2}}
\newcommand{\mpar}[3]{\frac{\partial^2 #1}{\partial #2 \partial #3}}  
\def\qed{ \hfill $\blacksquare$}  
%% End Local macros
%%%%%%%%%%%%%%%%%%%%%%%%%%%%%%%%%%%%%%%%%%%
%%%%%%%%%%%%%%%%%%%%%%%%%%%%%%%%%%%%%%%%%%%
%% Greek symbols
\let\ga=\alpha \let\gb=\beta \let\gc=\gamma \let\gd=\delta \let\gee=\epsilon
\let\gf=\varphi \let\gh=\eta \let\gi=\iota  \let\gk=\kappa \let\gl=\lambda \let\gm=\mu      \let\gn=\nu \let\go=\omega \let\gp=\pi \let\gr=\rho \let\gs=\sigma \let\gt=\tau \let\gth=\vartheta
\let\gx=\chi \let\gy=\upsilon \let\gz=\zeta
\let\gC=\Gamma \let\gD=\Delta \let\gF=\Phi \let\gL=\Lambda \let\gTh=\Theta
\let\gO=\Omega   \let\gP=\Pi    \let\gPs=\Psi  \let\gS=\Sigma \let\gU=\Upsilon \let\gX=\Chi
\let\gY=\Upsilon                                
%% End Greek Symbols
%%%%%%%%%%%%%%%%%%%%%%%%%%%%%%%%%%%%%%%%%%%
%%%%%%%%%%%%%%%%%%%%%%%%%%%%%%%%%%%%%%%%%%%
%% MathCal symbols
\newcommand{\mc}[1]{\mathcal{#1}}
\newcommand{\cA}{\mathcal{A}}\newcommand{\cB}{\mathcal{B}}\newcommand{\cC}{\mathcal{C}}
\newcommand{\cD}{\mathcal{D}}\newcommand{\cE}{\mathcal{E}}\newcommand{\cF}{\mathcal{F}}
\newcommand{\cG}{\mathcal{G}}\newcommand{\cH}{\mathcal{H}}\newcommand{\cI}{\mathcal{I}}
\newcommand{\cJ}{\mathcal{J}}\newcommand{\cK}{\mathcal{K}}\newcommand{\cL}{\mathcal{L}}
\newcommand{\cM}{\mathcal{M}}\newcommand{\cN}{\mathcal{N}}\newcommand{\cO}{\mathcal{O}}
\newcommand{\cP}{\mathcal{P}}\newcommand{\cQ}{\mathcal{Q}}\newcommand{\cR}{\mathcal{R}}
\newcommand{\cS}{\mathcal{S}}\newcommand{\cT}{\mathcal{T}}\newcommand{\cU}{\mathcal{U}}
\newcommand{\cV}{\mathcal{V}}\newcommand{\cW}{\mathcal{W}}\newcommand{\cX}{\mathcal{X}}
\newcommand{\cY}{\mathcal{Y}}\newcommand{\cZ}{\mathcal{Z}}  
%% End MathCal symbols
%%%%%%%%%%%%%%%%%%%%%%%%%%%%%%%%%%%%%%%%%%%
%%%%%%%%%%%%%%%%%%%%%%%%%%%%%%%%%%%%%%%%%%%
%% Math Boldface Symbols
\newcommand{\vzero}{\mathbf{0}}\newcommand{\vone}{\mathbf{1}}\newcommand{\vtwo}{\mathbf{2}}
\newcommand{\vthree}{\mathbf{3}}\newcommand{\vfour}{\mathbf{4}}\newcommand{\vfive}{\mathbf{5}}
\newcommand{\vsix}{\mathbf{6}}\newcommand{\vseven}{\mathbf{7}}\newcommand{\veight}{\mathbf{8}}
\newcommand{\vnine}{\mathbf{9}}\newcommand{\vA}{\mathbf{A}}\newcommand{\vB}{\mathbf{B}}
\newcommand{\vC}{\mathbf{C}}\newcommand{\vD}{\mathbf{D}}\newcommand{\vE}{\mathbf{E}}
\newcommand{\vF}{\mathbf{F}}\newcommand{\vG}{\mathbf{G}}\newcommand{\vH}{\mathbf{H}}
\newcommand{\vI}{\mathbf{I}}\newcommand{\vJ}{\mathbf{J}}\newcommand{\vK}{\mathbf{K}}
\newcommand{\vL}{\mathbf{L}}\newcommand{\vM}{\mathbf{M}}\newcommand{\vN}{\mathbf{N}}
\newcommand{\vO}{\mathbf{O}}\newcommand{\vP}{\mathbf{P}}\newcommand{\vQ}{\mathbf{Q}}
\newcommand{\vR}{\mathbf{R}}\newcommand{\vS}{\mathbf{S}}\newcommand{\vT}{\mathbf{T}}
\newcommand{\vU}{\mathbf{U}}\newcommand{\vV}{\mathbf{V}}\newcommand{\vW}{\mathbf{W}}
\newcommand{\vX}{\mathbf{X}}\newcommand{\vY}{\mathbf{Y}}\newcommand{\vZ}{\mathbf{Z}}
\newcommand{\va}{\mathbf{a}}\newcommand{\vb}{\mathbf{b}}\newcommand{\vc}{\mathbf{c}}
\newcommand{\vd}{\mathbf{d}}\newcommand{\ve}{\mathbf{e}}\newcommand{\vf}{\mathbf{f}}
\newcommand{\vg}{\mathbf{g}}\newcommand{\vh}{\mathbf{h}}\newcommand{\vi}{\mathbf{i}}
\newcommand{\vj}{\mathbf{j}}\newcommand{\vk}{\mathbf{k}}\newcommand{\vl}{\mathbf{l}}
\newcommand{\vm}{\mathbf{m}}\newcommand{\vn}{\mathbf{n}}\newcommand{\vo}{\mathbf{o}}
\newcommand{\vp}{\mathbf{p}}\newcommand{\vq}{\mathbf{q}}\newcommand{\vr}{\mathbf{r}}
\newcommand{\vs}{\mathbf{s}}\newcommand{\vt}{\mathbf{t}}\newcommand{\vu}{\mathbf{u}}
\newcommand{\vv}{\mathbf{v}}\newcommand{\vw}{\mathbf{w}}\newcommand{\vx}{\mathbf{x}}
\newcommand{\vy}{\mathbf{y}}\newcommand{\vz}{\mathbf{z}} 
%% End Math Boldface Symbols
%%%%%%%%%%%%%%%%%%%%%%%%%%%%%%%%%%%%%%%%%%%
%%%%%%%%%%%%%%%%%%%%%%%%%%%%%%%%%%%%%%%%%%%  
%% Math Bold Symbols commands  
\newcommand{\mv}[1]{\boldsymbol{#1}}\newcommand{\mvinfty}{\boldsymbol{\infty}}\newcommand{\mvzero}{\boldsymbol{0}}
\newcommand{\mvone}{\boldsymbol{1}}\newcommand{\mvtwo}{\boldsymbol{2}}\newcommand{\mvthree}{\boldsymbol{3}}
\newcommand{\mvfour}{\boldsymbol{4}}\newcommand{\mvfive}{\boldsymbol{5}}\newcommand{\mvsix}{\boldsymbol{6}}
\newcommand{\mvseven}{\boldsymbol{7}}\newcommand{\mveight}{\boldsymbol{8}}\newcommand{\mvnine}{\boldsymbol{9}}
\newcommand{\mvA}{\boldsymbol{A}}\newcommand{\mvB}{\boldsymbol{B}}\newcommand{\mvC}{\boldsymbol{C}}
\newcommand{\mvD}{\boldsymbol{D}}\newcommand{\mvE}{\boldsymbol{E}}\newcommand{\mvF}{\boldsymbol{F}}
\newcommand{\mvG}{\boldsymbol{G}}\newcommand{\mvH}{\boldsymbol{H}}\newcommand{\mvI}{\boldsymbol{I}}
\newcommand{\mvJ}{\boldsymbol{J}}\newcommand{\mvK}{\boldsymbol{K}}\newcommand{\mvL}{\boldsymbol{L}}
\newcommand{\mvM}{\boldsymbol{M}}\newcommand{\mvN}{\boldsymbol{N}}\newcommand{\mvO}{\boldsymbol{O}}
\newcommand{\mvP}{\boldsymbol{P}}\newcommand{\mvQ}{\boldsymbol{Q}}\newcommand{\mvR}{\boldsymbol{R}}
\newcommand{\mvS}{\boldsymbol{S}}\newcommand{\mvT}{\boldsymbol{T}}\newcommand{\mvU}{\boldsymbol{U}}
\newcommand{\mvV}{\boldsymbol{V}}\newcommand{\mvW}{\boldsymbol{W}}\newcommand{\mvX}{\boldsymbol{X}}
\newcommand{\mvY}{\boldsymbol{Y}}\newcommand{\mvZ}{\boldsymbol{Z}}\newcommand{\mva}{\boldsymbol{a}}
\newcommand{\mvb}{\boldsymbol{b}}\newcommand{\mvc}{\boldsymbol{c}}\newcommand{\mvd}{\boldsymbol{d}}
\newcommand{\mve}{\boldsymbol{e}}\newcommand{\mvf}{\boldsymbol{f}}\newcommand{\mvg}{\boldsymbol{g}}
\newcommand{\mvh}{\boldsymbol{h}}\newcommand{\mvi}{\boldsymbol{i}}\newcommand{\mvj}{\boldsymbol{j}}
\newcommand{\mvk}{\boldsymbol{k}}\newcommand{\mvl}{\boldsymbol{l}}\newcommand{\mvm}{\boldsymbol{m}}
\newcommand{\mvn}{\boldsymbol{n}}\newcommand{\mvo}{\boldsymbol{o}}\newcommand{\mvp}{\boldsymbol{p}}
\newcommand{\mvq}{\boldsymbol{q}}\newcommand{\mvr}{\boldsymbol{r}}\newcommand{\mvs}{\boldsymbol{s}}
\newcommand{\mvt}{\boldsymbol{t}}\newcommand{\mvu}{\boldsymbol{u}}\newcommand{\mvv}{\boldsymbol{v}}
\newcommand{\mvw}{\boldsymbol{w}}\newcommand{\mvx}{\boldsymbol{x}}\newcommand{\mvy}{\boldsymbol{y}}
\newcommand{\mvz}{\boldsymbol{z}}
\newcommand{\mvga}{\boldsymbol{\alpha}}\newcommand{\mvgb}{\boldsymbol{\beta}}\newcommand{\mvgc}{\boldsymbol{\gamma}}
\newcommand{\mvgC}{\boldsymbol{\Gamma}}\newcommand{\mvgd}{\boldsymbol{\delta}}\newcommand{\mvgD}{\boldsymbol{\Delta}}
\newcommand{\mvgee}{\boldsymbol{\epsilon}}\newcommand{\mveps}{\boldsymbol{\eps}}\newcommand{\mvgf}{\boldsymbol{\varphi}}
\newcommand{\mvgF}{\boldsymbol{\Phi}}\newcommand{\mvgth}{\boldsymbol{\theta}}\newcommand{\mvgTh}{\boldsymbol{\Theta}}
\newcommand{\mvgh}{\boldsymbol{\eta}}\newcommand{\mvgi}{\boldsymbol{\iota}}\newcommand{\mvgk}{\boldsymbol{\kappa}}
\newcommand{\mvgl}{\boldsymbol{\lambda}}\newcommand{\mvgL}{\boldsymbol{\Lambda}}\newcommand{\mvgm}{\boldsymbol{\mu}}
\newcommand{\mvgn}{\boldsymbol{\nu}}\newcommand{\mvgo}{\boldsymbol{\omega}}\newcommand{\mvgO}{\boldsymbol{\Omega}}
\newcommand{\mvgp}{\boldsymbol{\pi}}\newcommand{\mvgP}{\boldsymbol{\Pi}}\newcommand{\mvgr}{\boldsymbol{\rho}}
\newcommand{\mvgs}{\boldsymbol{\sigma}}\newcommand{\mvgS}{\boldsymbol{\Sigma}}\newcommand{\mvgt}{\boldsymbol{\tau}}
\newcommand{\mvgu}{\boldsymbol{\upsilon}}\newcommand{\mvgU}{\boldsymbol{\Upsilon}}\newcommand{\mvgx}{\boldsymbol{\chi}}
\newcommand{\mvgz}{\boldsymbol{\zeta}}\newcommand{\mvalpha}{\boldsymbol{\alpha}}\newcommand{\mvbeta}{\boldsymbol{\beta}}
\newcommand{\mvgamma}{\boldsymbol{\gamma}}\newcommand{\mvGamma}{\boldsymbol{\Gamma}}\newcommand{\mvdelta}{\boldsymbol{\delta}}
\newcommand{\mvDelta}{\boldsymbol{\Delta}}\newcommand{\mvepsilon}{\boldsymbol{\epsilon}}\newcommand{\mvphi}{\boldsymbol{\phi}}
\newcommand{\mvPhi}{\boldsymbol{\Phi}}\newcommand{\mvtheta}{\boldsymbol{\theta}}\newcommand{\mvTheta}{\boldsymbol{\Theta}}
\newcommand{\mveta}{\boldsymbol{\eta}}\newcommand{\mviota}{\boldsymbol{\iota}}\newcommand{\mvkappa}{\boldsymbol{\kappa}}
\newcommand{\mvlambda}{\boldsymbol{\lambda}}\newcommand{\mvLambda}{\boldsymbol{\Lambda}}\newcommand{\mvmu}{\boldsymbol{\mu}}
\newcommand{\mvnu}{\boldsymbol{\nu}}
\newcommand{\mvomega}{\boldsymbol{\omega}}\newcommand{\mvOmega}{\boldsymbol{\Omega}}\newcommand{\mvpi}{\boldsymbol{\pi}}
\newcommand{\mvPi}{\boldsymbol{\Pi}}\newcommand{\mvrho}{\boldsymbol{\rho}}\newcommand{\mvsigma}{\boldsymbol{\sigma}}
\newcommand{\mvSigma}{\boldsymbol{\Sigma}}\newcommand{\mvtau}{\boldsymbol{\tau}}\newcommand{\mvupsilon}{\boldsymbol{\upsilon}}
\newcommand{\mvUpsilon}{\boldsymbol{\Upsilon}}\newcommand{\mvchi}{\boldsymbol{\chi}}\newcommand{\mvzeta}{\boldsymbol{\zeta}}
\newcommand{\mvvartheta}{\boldsymbol{\vartheta}}\newcommand{\mvxi}{\boldsymbol{\xi}}\newcommand{\mvXi}{\boldsymbol{\Xi}}
\newcommand{\mvomic}{\boldsymbol{\o}}\newcommand{\mvvarpi}{\boldsymbol{\varpi}}\newcommand{\mvvarrho}{\boldsymbol{\varrho}}
\newcommand{\mvvarsigma}{\boldsymbol{\varsigma}}\newcommand{\mvvarphi}{\boldsymbol{\varphi}}\newcommand{\mvpsi}{\boldsymbol{\psi}}\newcommand{\mvPsi}{\boldsymbol{\Psi}}   
%% Math Bold Symbols commands  
%%%%%%%%%%%%%%%%%%%%%%%%%%%%%%%%%%%%%%%%%%%
%%%%%%%%%%%%%%%%%%%%%%%%%%%%%%%%%%%%%%%%%%%
%% Math Frakur fonts except for i which is \mfi
\newcommand{\f}[1]{\mathfrak{#1}}
\newcommand{\fA}{\mathfrak{A}}\newcommand{\fB}{\mathfrak{B}}\newcommand{\fC}{\mathfrak{C}}
\newcommand{\fD}{\mathfrak{D}}\newcommand{\fE}{\mathfrak{E}}\newcommand{\fF}{\mathfrak{F}}
\newcommand{\fG}{\mathfrak{G}}\newcommand{\fH}{\mathfrak{H}}\newcommand{\fI}{\mathfrak{I}}
\newcommand{\fJ}{\mathfrak{J}}\newcommand{\fK}{\mathfrak{K}}\newcommand{\fL}{\mathfrak{L}}
\newcommand{\fM}{\mathfrak{M}}\newcommand{\fN}{\mathfrak{N}}\newcommand{\fO}{\mathfrak{O}}
\newcommand{\fP}{\mathfrak{P}}\newcommand{\fQ}{\mathfrak{Q}}\newcommand{\fR}{\mathfrak{R}}
\newcommand{\fS}{\mathfrak{S}}\newcommand{\fT}{\mathfrak{T}}\newcommand{\fU}{\mathfrak{U}}
\newcommand{\fV}{\mathfrak{V}}\newcommand{\fW}{\mathfrak{W}}\newcommand{\fX}{\mathfrak{X}}
\newcommand{\fY}{\mathfrak{Y}}\newcommand{\fZ}{\mathfrak{Z}}\newcommand{\fa}{\mathfrak{a}}
\newcommand{\fb}{\mathfrak{b}}\newcommand{\fc}{\mathfrak{c}}\newcommand{\fd}{\mathfrak{d}}
\newcommand{\fe}{\mathfrak{e}}\newcommand{\ff}{\mathfrak{f}}\newcommand{\fg}{\mathfrak{g}}
\newcommand{\fh}{\mathfrak{h}}\newcommand{\mfi}{\mathfrak{i}}\newcommand{\fj}{\mathfrak{j}}
\newcommand{\fk}{\mathfrak{k}}\newcommand{\fl}{\mathfrak{l}}\newcommand{\fm}{\mathfrak{m}}
\newcommand{\fn}{\mathfrak{n}}\newcommand{\fo}{\mathfrak{o}}\newcommand{\fp}{\mathfrak{p}}
\newcommand{\fq}{\mathfrak{q}}\newcommand{\fr}{\mathfrak{r}}\newcommand{\fs}{\mathfrak{s}}
\newcommand{\ft}{\mathfrak{t}}\newcommand{\fu}{\mathfrak{u}}\newcommand{\fv}{\mathfrak{v}}
\newcommand{\fw}{\mathfrak{w}}\newcommand{\fx}{\mathfrak{x}}\newcommand{\fy}{\mathfrak{y}}
\newcommand{\fz}{\mathfrak{z}}  
%% End Math Frak                             
%%%%%%%%%%%%%%%%%%%%%%%%%%%%%%%%%%%%%%%%%%% 
%%%%%%%%%%%%%%%%%%%%%%%%%%%%%%%%%%%%%%%%%%%
% Double capital letters
\newcommand{\bb}[1]{\mathbb{#1}}
\newcommand{\bA}{\mathbb{A}}\newcommand{\bB}{\mathbb{B}}\newcommand{\bC}{\mathbb{C}}
\newcommand{\bD}{\mathbb{D}}\newcommand{\bE}{\mathbb{E}}\newcommand{\bF}{\mathbb{F}}
\newcommand{\bG}{\mathbb{G}}\newcommand{\bH}{\mathbb{H}}\newcommand{\bI}{\mathbb{I}}
\newcommand{\bJ}{\mathbb{J}}\newcommand{\bK}{\mathbb{K}}\newcommand{\bL}{\mathbb{L}}
\newcommand{\bM}{\mathbb{M}}\newcommand{\bN}{\mathbb{N}}\newcommand{\bO}{\mathbb{O}}
\newcommand{\bP}{\mathbb{P}}\newcommand{\bQ}{\mathbb{Q}}\newcommand{\bR}{\mathbb{R}}
\newcommand{\bS}{\mathbb{S}}\newcommand{\bT}{\mathbb{T}}\newcommand{\bU}{\mathbb{U}}
\newcommand{\bV}{\mathbb{V}}\newcommand{\bW}{\mathbb{W}}\newcommand{\bX}{\mathbb{X}}
\newcommand{\bY}{\mathbb{Y}}\newcommand{\bZ}{\mathbb{Z}}        
% End Double capital letters      
%%%%%%%%%%%%%%%%%%%%%%%%%%%%%%%%%%%%%%%%%%%
%%%%%%%%%%%%%%%%%%%%%%%%%%%%%%%%%%%%%%%%%%% 
%% Blackboard bold
\newcommand{\dA}{\mathds{A}}\newcommand{\dB}{\mathds{B}}
\newcommand{\dC}{\mathds{C}}\newcommand{\dD}{\mathds{D}}
\newcommand{\dE}{\mathds{E}}\newcommand{\dF}{\mathds{F}}
\newcommand{\dG}{\mathds{G}}\newcommand{\dH}{\mathds{H}}
\newcommand{\dI}{\mathds{I}}\newcommand{\dJ}{\mathds{J}}
\newcommand{\dK}{\mathds{K}}\newcommand{\dL}{\mathds{L}}
\newcommand{\dM}{\mathds{M}}\newcommand{\dN}{\mathds{N}}
\newcommand{\dO}{\mathds{O}}\newcommand{\dP}{\mathds{P}}
\newcommand{\dQ}{\mathds{Q}}\newcommand{\dR}{\mathds{R}}
\newcommand{\dS}{\mathds{S}}\newcommand{\dT}{\mathds{T}}
\newcommand{\dU}{\mathds{U}}\newcommand{\dV}{\mathds{V}}
\newcommand{\dW}{\mathds{W}}\newcommand{\dX}{\mathds{X}}
\newcommand{\dY}{\mathds{Y}}\newcommand{\dZ}{\mathds{Z}} 
%% Blackboard bold
%%%%%%%%%%%%%%%%%%%%%%%%%%%%%%%%%%%%%%%%%%%
%%%%%%%%%%%%%%%%%%%%%%%%%%%%%%%%%%%%%%%%%%%
% Roman capital letters
\newcommand{\rA}{\mathrm{A}}\newcommand{\rB}{\mathrm{B}}\newcommand{\rC}{\mathrm{C}}\newcommand{\rD}{\mathrm{D}}
\newcommand{\rE}{\mathrm{E}}\newcommand{\rF}{\mathrm{F}}\newcommand{\rG}{\mathrm{G}}\newcommand{\rH}{\mathrm{H}}
\newcommand{\rI}{\mathrm{I}}\newcommand{\rJ}{\mathrm{J}}\newcommand{\rK}{\mathrm{K}}\newcommand{\rL}{\mathrm{L}}
\newcommand{\rM}{\mathrm{M}}\newcommand{\rN}{\mathrm{N}}\newcommand{\rO}{\mathrm{O}}\newcommand{\rP}{\mathrm{P}}
\newcommand{\rQ}{\mathrm{Q}}\newcommand{\rR}{\mathrm{R}}\newcommand{\rS}{\mathrm{S}}\newcommand{\rT}{\mathrm{T}}
\newcommand{\rU}{\mathrm{U}}\newcommand{\rV}{\mathrm{V}}\newcommand{\rW}{\mathrm{W}}\newcommand{\rX}{\mathrm{X}}
\newcommand{\rY}{\mathrm{Y}}\newcommand{\rZ}{\mathrm{Z}}
% End Roman capital letters   
\newcommand{\sA}{\mathscr{A}}
\newcommand{\sB}{\mathscr{B}}
\newcommand{\sC}{\mathscr{C}}
\newcommand{\sD}{\mathscr{D}}
\newcommand{\sE}{\mathscr{E}}
\newcommand{\sF}{\mathscr{F}}
\newcommand{\sG}{\mathscr{G}}
\newcommand{\sH}{\mathscr{H}}
\newcommand{\sI}{\mathscr{I}}
\newcommand{\sJ}{\mathscr{J}}
\newcommand{\sK}{\mathscr{K}}
\newcommand{\sL}{\mathscr{L}}
\newcommand{\sM}{\mathscr{M}}
\newcommand{\sN}{\mathscr{N}}
\newcommand{\sO}{\mathscr{O}}
\newcommand{\sP}{\mathscr{P}}
\newcommand{\sQ}{\mathscr{Q}}
\newcommand{\sR}{\mathscr{R}}
\newcommand{\sS}{\mathscr{S}}
\newcommand{\sT}{\mathscr{T}}
\newcommand{\sU}{\mathscr{U}}
\newcommand{\sV}{\mathscr{V}}
\newcommand{\sW}{\mathscr{W}}
\newcommand{\sX}{\mathscr{X}}
\newcommand{\sY}{\mathscr{Y}}
\newcommand{\sZ}{\mathscr{Z}}




\newcommand{\subc}{\emph{sub-critical}}
\newcommand{\crit}{\emph{critical}}
\newcommand{\supc}{\emph{super-critical}}


%%%%%%%%%%%%%%%%%%%%%%%%%%%%%%%%%%%%%%%%%%%
%%%%%%%%%%%%%%%%%%%%%%%%%%%%%%%%%%%%%%%%%%%
%% Local Math Operator 
\DeclareMathOperator{\E}{\mathds{E}}
\DeclareMathOperator{\pr}{\mathds{P}}
\DeclareMathOperator{\sgn}{sgn}
\DeclareMathOperator{\var}{Var}
\DeclareMathOperator{\cov}{Cov}
\DeclareMathOperator{\argmax}{argmax}
\DeclareMathOperator{\argmin}{argmin}
\DeclareMathOperator{\hess}{Hess}
\DeclareMathOperator{\tr}{Tr} 
\DeclareMathOperator{\sech}{sech}  
\DeclareMathOperator{\diag}{diag}
\DeclareMathOperator{\fvec}{vec}
\DeclareMathOperator{\sym}{sym}
\DeclareMathOperator{\N}{N}
%% End Local Math Operator
%%%%%%%%%%%%%%%%%%%%%%%%%%%%%%%%%%%%%%%%%%%

\newcommand{\qiang}[1]{\textcolor{red}{#1}}

\newcommand{\ovR}{\overline{\vR}} 
\newcommand{\oR}{\overline{R}} 
\newcommand{\ov}{\overline} 
\newcommand{\indf}{\emph{indefinite}}
\newcommand{\pd}{\emph{positive-definite}}
\newcommand{\wmvgs}{\widetilde{\mvgs}}
\newcommand{\half}{\sfrac12}
\newcommand{\quar}{\sfrec14}

\newcommand{\hh}{\hat{h}}

\begin{document}

\title[Mixed $p$-spin free-energy fluctuations]{Free-energy fluctuations of mixed $p$-spin models throughout the high-temperature regime}
\author[Wu]{Qiang Wu}
\email{qiangw.math@gmail.com}
\date{}
% Compatibility with older amsart releases used for PDF export.
\makeatletter
\@ifundefined{subjclassname@2020}{\@namedef{subjclassname@2020}{\textup{2020} Mathematics Subject Classification}}{}
\makeatother
\subjclass[2020]{Primary: 82B26, 82B44, 60F05.}
\keywords{Spin glasses, high temperature, free-energy fluctuations, central limit theorem.}
%\thanks{}
%%%%%%%%%%%%%%%%%%%%%%%%%%%%%%%%%%%%%%%%%%%%%%%%%%%%%%%%%%%%%%%%%%%%%%%%%
\begin{abstract}
	We prove Gaussian fluctuations of the free energy for finite, zero-field, distinct-index mixed $p$-spin models whose minimum active interaction order is at least three. The result holds at every fixed inverse temperature strictly below the thermodynamic critical point. Uniformly on compact subintervals of this regime, the free energy admits a quadratic approximation in the Gaussian couplings with an explicit vanishing mean-square error at the fluctuation scale. The proof combines Chen's high-temperature overlap bounds with Gaussian interpolation and the Poincar\'e inequality. The limiting variance is determined by the smallest active interaction order, and the centering retains the exact finite-size covariance. 
\end{abstract}

%%%%%%%%%%%%%%%%%%%%%%%%%%%%%%%%%%%%%%%%%%%%%%%%%%%%%%%%%%%%%%%%%%%%%%%%%
\maketitle
\setcounter{tocdepth}{1}\tableofcontents


\section{Introduction and main results}
\label{sec:introduction}


The limiting free energy describes the typical thermodynamic behavior of
a spin glass, but does not determine the size or origin of its
sample-to-sample fluctuations. For interactions involving at least three
distinct spins, these fluctuations can be much smaller than the
finite-size corrections to the limiting free energy. A central question
is whether their leading contribution can nevertheless be described by
a simple statistic of the random couplings.

This paper studies that question for finite mixtures without an external
field. We identify the leading quadratic contribution throughout the
strict thermodynamic high-temperature regime and obtain a Gaussian limit
with the exact finite-size centering. The approximation also separates
two roles of the mixture: the smallest active interaction order determines
the leading random scale, whereas the whole covariance function determines
the high-temperature region.

\subsection{Setting}
\label{subsec:setting}
We first introduce the definition of the mixed $p$-spin models. Fix an integer $P\geq 3$ and real numbers
\(
\gamma_3,\ldots,\gamma_{P } \in \bR
\)
such that $\sum_{i=3}^{P} \abs{\gamma_i} \neq 0$.  Let
\begin{equation}
\cP
=
\left\{
p\in\{3,\ldots,P\}:\gamma_p\neq 0
\right\},
\qquad
p_\star
=
\min\cP.
\label{eq:active-degrees}
\end{equation}
For $N\geq P$, let
\(
\Sigma_N=\{-1,+1\}^N
\) be the configuration space.
For each $\sigma \in \Sigma_N$, the mixed $p$-spin Hamiltonian is 
\begin{equation}
H_N(\sigma)
=
\sum_{p\in\cP}
\frac{\gamma_p}{N^{(p-1)/2}}
\sum_{I\in\binom{[N]}{p}}
g_{p,I}\sigma_I,
\label{eq:Hamiltonian}
\end{equation}
where for every $p\in\cP$ and every $p$-hyperedge
\(
I\in\binom{[N]}{p},
\)
 $g_{p,I}$ are independent standard Gaussian random variables, and 
\(\sigma_I=\prod_{i\in I}\sigma_i.
\) For two configurations
$\sigma^1,\sigma^2\in\Sigma_N$, define their overlap by
\begin{equation}
R_{1,2}
=
\frac1N
\sum_{i=1}^N
\sigma_i^1\sigma_i^2.
\label{eq:overlap}
\end{equation}
For each $p\in\cP$, define the finite-$N$ distinct-index covariance
kernel by
\begin{equation}
\kappa_{p,N}(R_{1,2})
=
\frac1{N^p}
\sum_{I\in\binom{[N]}p}
\sigma_I^1\sigma_I^2.
\label{eq:kappa-pN}
\end{equation}
The quantity on the right-hand side depends on
$(\sigma^1,\sigma^2)$ only through $R_{1,2}$. A simple calculation gives
\begin{equation}
\E H_N(\sigma^1)H_N(\sigma^2)
=
N\xi_N(R_{1,2}).
\label{eq:covariance}
\end{equation}
where
\begin{equation}
\xi_N(q)
:=
\sum_{p\in\cP}
\gamma_p^2\kappa_{p,N}(q).
\label{eq:xi-N}
\end{equation}
In particular,
\begin{equation}
\xi_N(1)
=
\sum_{p\in\cP}
\gamma_p^2
\frac{\binom Np}{N^p}.
\label{eq:xi-N-one}
\end{equation}

The limiting covariance function is
\begin{equation}
\xi(q)
=
\sum_{p\in\cP}
\frac{\gamma_p^2}{p!}q^p,
\qquad
q\in[-1,1].
\label{eq:xi}
\end{equation}
In particular,
\begin{equation}
\xi(1)
=
\sum_{p\in\cP}
\frac{\gamma_p^2}{p!}.
\label{eq:xi-one}
\end{equation}

For $\beta\geq 0$, define the normalized partition function and the
free energy respectively
\begin{equation}
Z_N(\beta)
=
2^{-N}
\sum_{\sigma\in\Sigma_N}
\exp\left\{
\beta H_N(\sigma)
\right\},
\qquad
F_N(\beta)
=
\frac1N\log Z_N(\beta).
\label{eq:partition-free-energy}
\end{equation}
Since
\[
\E H_N(\sigma)^2=N\xi_N(1)
\]
for every $\sigma\in\Sigma_N$, one has
\begin{equation}
\frac1N\log\E Z_N(\beta)
=
\frac{\beta^2}{2}\xi_N(1).
\label{eq:annealed-free-energy}
\end{equation}

Let
\begin{equation}
F(\beta)
=
\lim_{N\to\infty}\E F_N(\beta)
\label{eq:limiting-free-energy}
\end{equation}
denote the limiting free energy.  We define the critical inverse
temperature by
\begin{equation}
\beta_c(\xi)
=
\sup\left\{
b\geq 0:
F(\beta)
=
\frac{\beta^2}{2}\xi(1)
\text{ for every }0\leq\beta\leq b
\right\}.
\label{eq:critical-inverse-temperature}
\end{equation}
Thus the strict high-temperature regime considered in this paper is
\begin{equation}
0\leq\beta<\beta_c(\xi).
\label{eq:high-temperature-regime}
\end{equation}
All estimates below are uniform on compact intervals
\begin{equation}
0\leq\beta\leq b,
\qquad
b<\beta_c(\xi).
\label{eq:compact-high-temperature-interval}
\end{equation}
We do not include the endpoint $\beta=\beta_c(\xi)$.

Define the quadratic disorder statistic
\begin{equation}
J_N(\beta)
=
\frac{\beta^2}{2}
\sum_{p\in\cP}
\frac{\gamma_p^2}{N^p}
\sum_{I\in\binom{[N]}p}
g_{p,I}^2.
\label{eq:J-N}
\end{equation}
Its expectation is
\begin{equation}
\E J_N(\beta)
=
\frac{\beta^2}{2}\xi_N(1).
\label{eq:mean-J-N}
\end{equation}


\subsection{Main results}
\label{subsec:main-results}

Our first result identifies the quadratic disorder statistic
$J_N(\beta)$ as the complete leading random contribution to the free
energy throughout the strict high-temperature regime.

\begin{thm}[Quadratic approximation]
\label{thm:quadratic-approximation}
Fix $0\leq b<\beta_c(\xi)$. There exists a constant
$C=C(b,\xi)<\infty$ such that, for every $N\geq P$,
\begin{equation}
\sup_{0\leq\beta\leq b}
\E
\left[
\left|
N^{p_\star/2}
\left(
F_N(\beta)-J_N(\beta)
\right)
\right|^2
\right]
\leq
C N^{1-p_\star/2}.
\label{eq:quadratic-approximation-rate}
\end{equation}
Consequently,
\begin{equation}
\lim_{N\to\infty}
\sup_{0\leq\beta\leq b}
\E
\left[
\left|
N^{p_\star/2}
\left(
F_N(\beta)-J_N(\beta)
\right)
\right|^2
\right]
=
0.
\label{eq:quadratic-approximation-limit}
\end{equation}
\end{thm}

The second result gives the leading free-energy fluctuation.  The
fluctuation scale and limiting variance are determined by the smallest
active interaction degree $p_\star$.

\begin{thm}[CLT of free energy]
\label{thm:main-clt}
For every fixed
\[
0<\beta<\beta_c(\xi),
\]
one has
\begin{equation}
N^{p_\star/2}
\left(
F_N(\beta)
-
\frac{\beta^2}{2}\xi_N(1)
\right)
\Longrightarrow
\mathcal N
\left(
0,
\frac{\beta^4\gamma_{p_\star}^4}{2p_\star!}
\right).
\label{eq:main-clt}
\end{equation}
\end{thm}

The finite-$N$ centering in Theorem~\ref{thm:main-clt} is essential.
Indeed, expanding the falling factorial gives
\begin{equation}
\xi_N(1)-\xi(1)
=-\frac1N\sum_{p\in\cP}\frac{\gamma_p^2}{p!}\binom p2+O(N^{-2}),
\label{eq:centering-difference}
\end{equation}
whereas the random fluctuation scale in
\eqref{eq:main-clt} is $N^{-p_\star/2}$.  Since $p_\star\geq 3$, the
coefficient of $N^{-1}$ in~\eqref{eq:centering-difference} is strictly negative for a nontrivial mixture. Thus the deterministic centering error has order $N^{-1}$, which is larger than the fluctuation scale for each fixed $\beta>0$.



\subsection{Related work and motivation}
\label{subsec:related-work}

The thermodynamic limit and the fluctuations around it are different
questions. Talagrand's proof of the Parisi formula~\cite{Talagrand2006}
and Panchenko's extension to mixed models with odd
interactions~\cite{Panchenko2014} provide the variational description of
the limiting free energy. They do not, by themselves, specify a central
limit scale or the finite-size centering at that scale. The definition
of $\beta_c(\xi)$ above uses this thermodynamic limit. A main issue in
fluctuation theory is whether estimates obtained in an explicit
sufficient high-temperature range continue to hold throughout the
region where the quenched and annealed limits agree.

For the zero-field Sherrington--Kirkpatrick model, Aizenman, Lebowitz,
and Ruelle~\cite{ALR1987} established high-temperature Gaussian
fluctuations through a graphical expansion whose leading terms are
loops. Comets and Neveu~\cite{CometsNeveu1995} developed a different
approach using Brownian motions and continuous martingales. These two
methods lead to distinct ways of studying higher-order interactions:
one identifies the combinatorial structures contributing to the
partition function, while the other controls a stochastic
representation of its fluctuations. Both approaches require estimates
that are sharper than concentration of the free energy around its
thermodynamic limit.

For pure distinct-index models with $p\geq3$, Bovier, Kurkova, and
L\"owe~\cite{BovierKurkovaLowe2002} combined martingale arguments with
truncation to obtain a Gaussian limit for
$\log(Z_N/\E Z_N)$ on the scale $N^{-(p-2)/2}$ in an explicit
high-temperature range. The corresponding scale for the free energy
per spin is $N^{-p/2}$. Kn\"opfel and L\"owe~\cite{KnopfelLowe2005}
extended the pure-model conclusion to symmetric, variance-one
couplings with an exponential moment, again in a specified
high-temperature range. Their factorization into a product of
hyperbolic cosines and a remaining spin average shows why the
fluctuations of squared couplings enter. Thus both the pure-model
Gaussian limit and its leading quadratic mechanism have direct
precedents.

The temperature restriction in these results deserves separate
attention. Talagrand~\cite{Talagrand2000} used a truncated
second-moment argument to obtain an explicit lower bound on the
critical inverse temperature of the pure model. Bovier and
Schertzer~\cite{BovierSchertzer2024} subsequently established the
pure-model fluctuation theorem throughout the range given by that
bound. They also identify the quadratic coupling statistic as the
complete leading contribution and prove a further central limit
theorem after subtracting it. Their paper explicitly distinguishes
this threshold from the actual critical inverse temperature. It uses
a normalization in which the Hamiltonian variance is exactly $N$.
The comparison with the present theorem therefore involves the
temperature range, the treatment of mixtures, and the strength of the
approximation, as well as the finite-size convention.

Mixed models with a two-spin component have a different leading
structure. Banerjee and Belius~\cite{BanerjeeBelius2021} prove a
Gaussian limit at sufficiently high temperature for canonical mixed
models with covariance $N\xi(R_{1,2})$ and a nonzero two-spin
coefficient. The order-one fluctuations of $\log Z_N$ arise from
weighted cycles formed from that component. Consequently, the
assumption $p_\star\geq3$ in the present paper separates two different
fluctuation mechanisms; it is not a restriction that can be removed
by inserting a two-spin term into the variance formula.

The closest graphical treatment of higher-order mixtures is due to
Dey and Wu~\cite{DeyWuHypergraphs}. They study distinct-index mixed
models with symmetric disorder under moment assumptions and with weak
external fields. Their decomposition separates the product of
hyperbolic cosines from an interacting factor, whose
fluctuations are described by leading hypergraphs. At zero field,
the relevant structures depend on the smallest active even and odd
orders. Their fluctuation analysis works below the explicit
second-moment threshold
\begin{equation*}
\beta_f^2=\inf_{0<q<1}\frac{I(q)}{\xi(q)},
\qquad
I(q)=\frac{(1+q)\log(1+q)+(1-q)\log(1-q)}2.
\end{equation*}
This threshold is not identified with $\beta_c(\xi)$. Their results
resolve the interacting contribution and allow disorder and fields beyond
the Gaussian zero-field setting considered here.

A second line of work studies fluctuations through overlap estimates.
For the pure $p$-spin model with a field, Bardina,
M\'arquez-Carreras, Rovira, and Tindel~\cite{BardinaEtAl2004} developed
cavity estimates that bound centered overlap moments at their natural
orders when the inverse temperature is sufficiently small. Such
estimates give quantitative control of replica-symmetric behavior,
but their temperature hypothesis is stronger than merely lying below
the thermodynamic critical point. For the zero-field higher-order
models relevant here, Chen~\cite{ChenHighTemperature} proves overlap
moment bounds throughout the strict thermodynamic high-temperature
region; Theorem~4 and the mixed-model extension in Remark~3 are the
specific inputs used below.

Chen's theorem is stated for the canonical covariance, so its
application to \eqref{eq:Hamiltonian} requires a comparison. After
matching the normalization, we transfer its bounds to the
distinct-index model by replicated Gaussian interpolation. The
remaining argument combines cancellations in the finite-index
covariance kernels with Gaussian integration by parts and the
Poincar\'e inequality. This gives a mean-square estimate for the
error after removing the quadratic statistic. The central limit
theorem then follows from an independent sum, rather than from a
classification of the leading interacting hypergraphs.

External fields illustrate why one must specify the regime before
comparing fluctuation theorems. Guerra and
Toninelli~\cite{GuerraToninelli2002} obtained Gaussian fluctuations
for the SK model with a nonzero field in a sufficiently
high-temperature region using quadratic replica coupling and cavity
equations. Tindel~\cite{Tindel2005} gave a stochastic-calculus
treatment in a sufficiently high-temperature region. For mixed
models, Chen, Dey, and Panchenko~\cite{ChenDeyPanchenko2017} prove
variance of order $N$ for the extensive free energy when the field is
nonzero; for mixtures without odd interaction terms of order at least
three, they also prove a central limit theorem at every temperature.
The resulting $N^{-1/2}$ scale for the free energy per spin differs
from the zero-field scale in Theorem~\ref{thm:main-clt}.

The passage from zero field to a fixed nonzero field is itself a
fluctuation problem. Dey and Wu~\cite{DeyWuWeakField2023} study SK
models with $h_N=\rho N^{-\alpha}$ and identify a transition at
$\alpha=1/4$: weaker fields leave the leading zero-field law
unchanged, the boundary scale changes its parameters, and stronger
vanishing fields produce a larger fluctuation scale. For
$\alpha\geq1/4$, their results cover every fixed $\beta<1$ and
symmetric disorder with a finite sixth moment. For $\alpha<1/4$,
the stated theorem assumes Gaussian disorder and sufficiently small
$\beta$. Their expansion introduces paths in addition to loops.
These results explain why a theorem at zero field does not
automatically describe a sequence of small nonzero fields.

Inhomogeneity across species supplies another useful comparison.
Dey and Wu~\cite{DeyWuMultiSpecies2021} develop a species-wise cavity
method for the multispecies SK model, including indefinite
interaction matrices. They establish overlap concentration and a
positive-field free-energy central limit theorem in a specified
high-temperature region; their zero-field overlap bound extends to
the spectral threshold. Their weak-field
work~\cite{DeyWuWeakField2023} also treats multispecies
fluctuations, with the bipartite case developed in detail.
Species-dependent two-spin variances lead to vector overlaps.
They differ from a mixture of interaction orders acting on the same
population of spins, for which the scalar covariance $\xi$ used
here is the relevant object.

Two further distinctions concern the strength and the uniformity of
fluctuation estimates. Chatterjee's
superconcentration theorem~\cite{Chatterjee2009} gives a
sublinear variance bound for the extensive zero-field SK free
energy at every temperature; such a bound alone does not identify
its precise limiting law. At the high-temperature boundary,
Dey and Kang~\cite{DeyKang2026} prove a variance asymptotic and a
Gaussian limit for SK inverse temperatures approaching $1$ with
$N^{1/3}(1-\beta_N^2)$ bounded away from zero. This requires
estimates sensitive to the distance from criticality. The
compact-interval uniformity in
Theorem~\ref{thm:quadratic-approximation} does not provide such an
$N$-dependent approach to $\beta_c(\xi)$.

Finally, thermodynamic equivalence is insufficient for transferring
these small fluctuation scales. Bovier, Kurkova, and
L\"owe~\cite{BovierKurkovaLowe2002} also analyze a repeated-index
convention and obtain a different scale. The exact-variance
normalization of~\cite{BovierSchertzer2024} likewise differs from
\eqref{eq:Hamiltonian} by an $N$-dependent factor. Here the
annealed center is exactly $\beta^2\xi_N(1)/2$; replacing it by
$\beta^2\xi(1)/2$ introduces the nonzero correction of order
$N^{-1}$ in \eqref{eq:centering-difference}, which is larger than
the random scale. Retaining that correction is necessary for the
stated limit.

The present result places the established quadratic mechanism in a
finite-mixture theorem with exact centering and a uniform
mean-square error bound throughout every fixed compact subinterval
of the strict thermodynamic high-temperature region. Only the
smallest active interaction order survives in the limiting
variance, whereas the whole mixture can affect $\beta_c(\xi)$.
The uniformity is the supremum, over temperature, of the
mean-square error in
\eqref{eq:quadratic-approximation-rate}. No convergence of a
temperature-indexed process, critical-endpoint limit, or
second-scale theorem after subtraction of $J_N$ is asserted.

\section{Proof of Theorem~\ref{thm:quadratic-approximation}}
\label{sec:quadratic-approximation}

The case $b=0$ is immediate. Fix $0<b<\beta_c(\xi)$. All constants in this section are chosen on the entire interval $0\leq s\leq b^2$ and may depend on $b$ and the fixed mixture, but not on $N$ or on the evaluation point $\beta_0\in[0,b]$. For $s\in[0,b^2]$, let $\langle\cdot\rangle_s$ denote the
Gibbs expectation associated with the Hamiltonian $\sqrt{s}\,H_N(\gs)$,
namely
\begin{equation}
G_{N,s}(\sigma)
=
\frac{2^{-N}\exp\{\sqrt{s}\,H_N(\sigma)\}}
{Z_N(\sqrt{s})}.
\label{eq:Gibbs-s}
\end{equation}
For a function $f: \Sigma_N^{\otimes \ell} \to \bR$ with $\ell$ finite, write
\begin{equation}
\nu_s(f)=\E\langle f\rangle_s.
\label{eq:nu-s}
\end{equation}
At $s=0$, the associated Gibbs measure is uniform on $\Sigma_N$. The
corresponding expectation will also be denoted by $\nu_0$. For an arbitrary $\beta_0\in[0,b]$, set
\begin{equation}
D_N
=
\log Z_N(\beta_0)-N J_N(\beta_0).
\label{eq:D-N-beta0}
\end{equation}
The scaled error in Theorem~\ref{thm:quadratic-approximation} is given by the identity 
\begin{equation}
N^{p_\star/2}
\bigl(F_N(\beta_0)-J_N(\beta_0)\bigr)
=
N^{p_\star/2-1}D_N.
\label{eq:scaled-D-N}
\end{equation}
We will control both the mean and the variance of $D_N$ through the
same quantity. Define
\begin{equation}
A_N(s)
=
\nu_s\bigl(\xi_N(R_{1,2})\bigr),
\qquad
0\leq s\leq b^2,
\label{eq:A-N-def}
\end{equation}
and
\begin{equation}
V_N
=
\nu_0\bigl(\xi_N(R_{1,2})^2\bigr)
=
\sum_{p\in\cP}
\gamma_p^4\frac{\binom Np}{N^{2p}}.
\label{eq:V-N-def}
\end{equation}
The second identity in \eqref{eq:V-N-def} follows from the
orthogonality of distinct Walsh characters under the uniform measure.
In particular,
\begin{equation}
V_N=O\bigl(N^{-p_\star}\bigr).
\label{eq:V-N-order}
\end{equation}
Also, using two independent Gibbs replicas,
\begin{equation}
A_N(s)
=
\sum_{p\in\cP}
\frac{\gamma_p^2}{N^p}
\sum_{I\in\binom{[N]}p}
\E\langle\sigma_I\rangle_s^2
\geq 0.
\label{eq:A-N-positive}
\end{equation}

The central estimate of the section is the following linear-response
statement.

\begin{prop}[Linear response of the covariance kernel]
\label{prop:linear-response}
For the fixed bound $b<\beta_c(\xi)$, one has
\begin{equation}
A_N(s)
=
sN V_N
+
O_{b,\xi}\bigl(N^{2-3p_\star/2}\bigr)
\label{eq:linear-response}
\end{equation}
uniformly for $0\leq s\leq b^2$.
\end{prop}

The need for the precise leading term in
\eqref{eq:linear-response} comes from the Gaussian Poincar\'e estimate
used at the end of the section. Its two principal terms are
\begin{equation}
\beta_0^2N A_N(\beta_0^2)-\beta_0^4N^2V_N.
\label{eq:cancellation-preview}
\end{equation}
Proposition~\ref{prop:linear-response} makes the leading parts of the
two terms in \eqref{eq:cancellation-preview} cancel exactly. The next
three subsections are devoted to the proof of
Proposition~\ref{prop:linear-response}. We return to $D_N$ and prove
Theorem~\ref{thm:quadratic-approximation} only after that proposition
has been established.

\subsection{Canonical comparison and overlap moments}
\label{subsec:canonical-overlap}

The proof of Proposition~\ref{prop:linear-response} requires a
microscopic moment bound for the overlap under the distinct-index
Gibbs measure. The available full-high-temperature estimate is
formulated for the canonical mixed $p$-spin model with covariance
$N\xi(R_{1,2})$. We therefore introduce that model only as an
auxiliary device for importing the overlap estimate. We do not use it
to compare the microscopic free-energy fluctuations of the two
models.

Let $\widetilde H_N$ be a centered Gaussian process on $\Sigma_N$,
independent of $H_N$, with covariance
\begin{equation}
\E\widetilde H_N(\sigma^1)\widetilde H_N(\sigma^2)
=
N\xi(R_{1,2}).
\label{eq:canonical-covariance}
\end{equation}
The following elementary comparison is the reason the canonical model
can be used.

\begin{lem}[Bounded covariance difference]
\label{lem:bounded-covariance-difference}
There exists $C_\xi<\infty$ such that
\begin{equation}
\sup_{\sigma^1,\sigma^2\in\Sigma_N}
\left|
N\xi_N(R_{1,2})-N\xi(R_{1,2})
\right|
\leq C_\xi
\label{eq:bounded-covariance-difference-proof}
\end{equation}
for every $N\geq P$.
\end{lem}

\begin{proof}
Fix $p\in\cP$ and put
\[
x_i=\sigma_i^1\sigma_i^2,
\qquad
r=\frac1N\sum_{i=1}^N x_i.
\]
Let $I_1,\ldots,I_p$ be independent and uniform on $[N]$, and let
$\mathcal D$ be the event that these indices are all distinct. Then
\[
\E\prod_{a=1}^p x_{I_a}=r^p,
\]
whereas
\[
\E\left[
\prod_{a=1}^p x_{I_a}\,\mathbf 1_{\mathcal D}
\right]
=
\frac{p!}{N^p}
\sum_{I\in\binom{[N]}p}\prod_{i\in I}x_i
=
p!\,\kappa_{p,N}(r).
\]
Since the product has absolute value one,
\begin{equation}
\left|
\kappa_{p,N}(r)-\frac{r^p}{p!}
\right|
\leq
\frac1{p!}\pr(\mathcal D^c)
\leq
\frac1{p!}\frac{\binom p2}{N}.
\label{eq:kappa-covariance-comparison}
\end{equation}
Multiplying by $N\gamma_p^2$, summing over the finite set $\cP$, and
using the definitions of $\xi_N$ and $\xi$ proves
\eqref{eq:bounded-covariance-difference-proof}.
\end{proof}

For later use, set
\begin{equation}
\Delta_N(\sigma,\tau)
=
N\xi_N(R(\sigma,\tau))-N\xi(R(\sigma,\tau)).
\label{eq:Delta-N}
\end{equation}
Then Lemma~\ref{lem:bounded-covariance-difference} states that
$\|\Delta_N\|_\infty\leq C_\xi$.

For completeness, Lemma~\ref{lem:bounded-covariance-difference} also
shows that the distinct-index and canonical models have the same
limiting free energy. Let
\[
\Phi_N(t)
=
\frac1N\E\log\left(
2^{-N}\sum_{\sigma\in\Sigma_N}
\exp\left\{
\beta\left(\sqrt{t}\,H_N(\sigma)
+\sqrt{1-t}\,\widetilde H_N(\sigma)\right)
\right\}
\right).
\]
Gaussian integration by parts gives
\[
\Phi_N'(t)
=
\frac{\beta^2}{2N}
\E\left\langle
\Delta_N(\sigma^1,\sigma^1)
-\Delta_N(\sigma^1,\sigma^2)
\right\rangle_t.
\]
Consequently, for every fixed $\beta$,
\begin{equation}
\left|
\E F_N(\beta)-\E\widetilde F_N(\beta)
\right|
\leq
\frac{\beta^2 C_\xi}{N}.
\label{eq:free-energy-comparison}
\end{equation}
Thus the critical inverse temperature defined in terms of the limiting
free energy is the same for both models.

We next record the overlap input from Theorem~4 and Remark~3 of~\cite{ChenHighTemperature}. To match conventions, put $a=\xi(1)>0$ and $c_p=\gamma_p^2/(p!a)$. Then $c_p\geq0$, $\sum_pc_p=1$, and Chen's normalized covariance is $\xi_C(q)=\tfrac12\sum_pc_pq^p=\xi(q)/(2a)$. The canonical model with covariance $N\xi$ at inverse temperature $\sqrt{s}$ therefore corresponds to the normalized model at inverse temperature $\sqrt{2as}$. The annealed free energies and critical inverse temperatures scale by the same change of variable. The thermodynamic limit for this canonical model is the background limit used in that result; the comparison above transfers it to the distinct-index model. The cited theorem and remark imply that for every integer
$k\geq1$ there exists $C=C(b,k,\xi)$ such that
\begin{equation}
\sup_{0\leq s\leq b^2}
\widetilde\nu_s\bigl(|R_{1,2}|^{2k}\bigr)
\leq
\frac{C}{N^k}.
\label{eq:canonical-overlap-moments}
\end{equation}
Here $\widetilde\nu_s$ denotes the disorder-averaged Gibbs expectation
for the canonical Hamiltonian $\sqrt{s}\,\widetilde H_N$. The
assumption $p_\star\geq3$ is exactly the lowest-degree hypothesis in
Remark~3 of \cite{ChenHighTemperature}.

We now transfer \eqref{eq:canonical-overlap-moments} to the
model of this paper.

\begin{lem}[Overlap moments for the distinct-index model]
\label{lem:distinct-overlap-moments}
For every integer $k\geq1$, there exists
$C=C(b,k,\xi)<\infty$ such that
\begin{equation}
\sup_{0\leq s\leq b^2}
\nu_s\bigl(|R_{1,2}|^{2k}\bigr)
\leq
\frac{C}{N^k}.
\label{eq:distinct-overlap-moments}
\end{equation}
\end{lem}

\begin{proof}
For $t\in[0,1]$, define the interpolating Hamiltonian
\begin{equation}
H_{N,t}(\sigma)
=
\sqrt{t}\,H_N (\sigma)+\sqrt{1-t}\,\widetilde H_N(\sigma),
\label{eq:canonical-interpolation}
\end{equation}
and let $\nu_{s,t}$ be the disorder-averaged replicated Gibbs
expectation associated with $\sqrt{s}\,H_{N,t}$. Recall from
\eqref{eq:Delta-N} that
\[
\|\Delta_N\|_\infty\leq C_\xi.
\]

We use the replicated Gaussian interpolation identity. If $f$ is a
bounded observable of $n$ replicas, then
\begin{align}
\frac{d}{dt}\nu_{s,t}(f)
=
\frac{s}{2}\nu_{s,t}\Bigg(
 f\Bigg[&
\sum_{\ell,\ell'=1}^n\Delta_{\ell,\ell'}
-2n\sum_{\ell=1}^n\Delta_{\ell,n+1}
\notag\\
&+n(n+1)\Delta_{n+1,n+2}
-n\Delta_{n+1,n+1}
\Bigg]
\Bigg),
\label{eq:replicated-interpolation}
\end{align}
where
\[
\Delta_{\ell,\ell'}
=
\Delta_N(\sigma^\ell,\sigma^{\ell'}).
\]
The last diagonal term in \eqref{eq:replicated-interpolation} is part
of the general identity. In the present model all diagonal values of
$\Delta_N$ are equal, so the diagonal terms can also be canceled and
\eqref{eq:replicated-interpolation} can be written in the usual
pair-overlap form. Formula \eqref{eq:replicated-interpolation} follows
by differentiating the replicated Gibbs expectation and applying
Gaussian integration by parts to each of the two Gaussian processes in
\eqref{eq:canonical-interpolation}.

Apply \eqref{eq:replicated-interpolation} with $n=2$ and
$f=|R_{1,2}|^{2k}$. Since $f\geq0$, $s\leq b^2$, and
$\|\Delta_N\|_\infty\leq C_\xi$, there exists
$C=C(b,k,\xi)$ such that
\begin{equation}
\left|
\frac{d}{dt}\nu_{s,t}(f)
\right|
\leq
C\nu_{s,t}(f).
\label{eq:gronwall-differential}
\end{equation}
Gronwall's inequality, followed by
\eqref{eq:canonical-overlap-moments}, gives
\[
\nu_{s,1}(f)
\leq
e^C\nu_{s,0}(f)
\leq
\frac{C'}{N^k}
\]
uniformly for $0\leq s\leq b^2$. Since $\nu_{s,1}=\nu_s$,
this proves \eqref{eq:distinct-overlap-moments}.
\end{proof}

\subsection{From overlap moments to covariance-kernel moments}
\label{subsec:kernel-moments}

The next lemma gives a finite-$N$ bound on the distinct-index kernel.

\begin{lem}[Finite-$N$ polynomial bound]
\label{lem:finite-kernel-bound}
For every fixed integer $p\geq1$, there exists $C_p<\infty$ such that,
for every $x_1,\ldots,x_N\in\{-1,+1\}$ and
$r=N^{-1}\sum_{i=1}^N x_i$,
\begin{equation}
\left|
\frac1{N^p}
\sum_{I\in\binom{[N]}p}\prod_{i\in I}x_i
\right|
\leq
C_p
\sum_{j=0}^{\lfloor p/2\rfloor}
N^{-j}|r|^{p-2j}.
\label{eq:finite-kernel-bound}
\end{equation}
\end{lem}

\begin{proof}
Let $e_p(x_1,\ldots,x_N)$ be the $p$th elementary symmetric
polynomial. For a partition $\lambda\vdash p$, write
$\ell(\lambda)$ for its number of parts, $o(\lambda)$ for the number of
odd parts, and
\[
z_\lambda
=
\prod_{m\geq1}m^{a_m(\lambda)}a_m(\lambda)!,
\]
where $a_m(\lambda)$ is the multiplicity of the part $m$. The standard
power-sum expansion of the elementary symmetric polynomial is
\begin{equation}
e_p
=
\sum_{\lambda\vdash p}
(-1)^{p-\ell(\lambda)}z_\lambda^{-1}
\prod_{a=1}^{\ell(\lambda)}
\left(\sum_{i=1}^N x_i^{\lambda_a}\right).
\label{eq:elementary-power-sum}
\end{equation}
Since $x_i\in\{-1,+1\}$,
\[
\sum_{i=1}^N x_i^m
=
\begin{cases}
Nr,&m\text{ odd},\\
N,&m\text{ even}.
\end{cases}
\]
Consequently, the absolute value of the term indexed by $\lambda$ in
$N^{-p}e_p$ is
\begin{equation}
z_\lambda^{-1}
N^{\ell(\lambda)-p}|r|^{o(\lambda)}.
\label{eq:partition-term}
\end{equation}
Set $j=(p-o(\lambda))/2$. The parity of a partition gives
$j\in\{0,\ldots,\lfloor p/2\rfloor\}$ and
$o(\lambda)=p-2j$. Moreover,
\begin{equation}
p-\ell(\lambda)
=
\sum_{a=1}^{\ell(\lambda)}(\lambda_a-1)
\geq
\frac{p-o(\lambda)}2
=j.
\label{eq:partition-count}
\end{equation}
Indeed, an odd part $m$ contributes $(m-1)/2$ to
$(p-o(\lambda))/2$ and an even part $m$ contributes $m/2$, while both
are at most $m-1$. Thus \eqref{eq:partition-term} is bounded by
$z_\lambda^{-1}N^{-j}|r|^{p-2j}$. Summing over the finitely many
partitions of $p$ proves \eqref{eq:finite-kernel-bound}.
\end{proof}

We now obtain the only kernel moment estimate needed below.

\begin{lem}[Cubic covariance-kernel moment]
\label{lem:cubic-kernel-moment}
There exists $C=C(b,\xi)<\infty$ such that
\begin{equation}
\sup_{0\leq s\leq b^2}
\nu_s\left(
\left|\xi_N(R_{1,2})\right|^3
\right)
\leq
C N^{-3p_\star/2}.
\label{eq:cubic-kernel-moment}
\end{equation}
\end{lem}

\begin{proof}
Fix $p\in\cP$. By Lemma~\ref{lem:finite-kernel-bound},
\[
|\kappa_{p,N}(R_{1,2})|
\leq
C_p\sum_{j=0}^{\lfloor p/2\rfloor}
N^{-j}|R_{1,2}|^{p-2j}.
\]
Choose an integer $k$ so large that $2k\geq3P$. By
Lemma~\ref{lem:distinct-overlap-moments},
\[
\sup_{0\leq s\leq b^2}
\|R_{1,2}\|_{L^{2k}(\nu_s)}
\leq
C N^{-1/2}.
\]
Since $3(p-2j)\leq2k$, monotonicity of $L^q$ norms gives
\begin{equation}
\sup_{0\leq s\leq b^2}
\left\|
N^{-j}R_{1,2}^{p-2j}
\right\|_{L^3(\nu_s)}
\leq
C N^{-j}N^{-(p-2j)/2}
=
C N^{-p/2}.
\label{eq:one-kernel-L3}
\end{equation}
The case $p-2j=0$ is included in the same bound. Hence
\[
\sup_{0\leq s\leq b^2}
\|\kappa_{p,N}(R_{1,2})\|_{L^3(\nu_s)}
\leq
C N^{-p/2}.
\]
Using the finite-mixture representation
$\xi_N=\sum_{p\in\cP}\gamma_p^2\kappa_{p,N}$ and the triangle
inequality in $L^3$, we obtain
\[
\sup_{0\leq s\leq b^2}
\|\xi_N(R_{1,2})\|_{L^3(\nu_s)}
\leq
C N^{-p_\star/2},
\]
which is equivalent to \eqref{eq:cubic-kernel-moment}.
\end{proof}

\subsection{Proof of the linear-response proposition}
\label{subsec:proof-linear-response}

We first record the stochastic-stability identity used in the proof.

\begin{lem}[Stochastic stability]
\label{lem:stochastic-stability}
Let $f$ be a bounded observable of $n$ replicas that does not depend on
the disorder. With
\[
\xi_{\ell,\ell'}
=
\xi_N(R_{\ell,\ell'}),
\]
one has, for $s>0$,
\begin{equation}
\frac{d}{ds}\nu_s(f)
=
N\nu_s\Bigg(
 f\Bigg[
\sum_{1\leq\ell<\ell'\leq n}\xi_{\ell,\ell'}
-n\sum_{\ell=1}^n\xi_{\ell,n+1}
+\frac{n(n+1)}2\xi_{n+1,n+2}
\Bigg]
\Bigg).
\label{eq:stochastic-stability}
\end{equation}
The identity extends to $s=0$ by continuity.
\end{lem}

\begin{proof}
Differentiate the replicated Gibbs expectation associated with
$\sqrt{s}\,H_N$ and apply Gaussian integration by parts. In the
general formula, the diagonal covariance terms are
\[
\frac12\left(
\sum_{\ell=1}^n N\xi_N(1)-nN\xi_N(1)
\right)=0,
\]
because the diagonal covariance is constant. The remaining
pair-overlap terms are exactly those in
\eqref{eq:stochastic-stability}. For each fixed $N$, the spin space is finite and the differentiated Gibbs observables have integrable Gaussian bounds. After integration by parts, the displayed right-hand side is continuous at $s=0$, which justifies the stated endpoint extension.
\end{proof}

\begin{proof}[Proof of Proposition~\ref{prop:linear-response}]
Since $A_N(s)=\nu_s(\xi_{1,2})$, applying
Lemma~\ref{lem:stochastic-stability} with $n=2$ gives
\begin{equation}
A_N'(s)
=
N\nu_s\left(
\xi_{1,2}^2
-4\xi_{1,2}\xi_{1,3}
+3\xi_{1,2}\xi_{3,4}
\right).
\label{eq:A-N-derivative}
\end{equation}
Introduce
\begin{align}
B_{1,N}(s)&=\nu_s(\xi_{1,2}^2),
\label{eq:B1-def}\\
B_{2,N}(s)&=\nu_s(\xi_{1,2}\xi_{1,3}),
\label{eq:B2-def}\\
B_{3,N}(s)&=\nu_s(\xi_{1,2}\xi_{3,4}).
\label{eq:B3-def}
\end{align}
At $s=0$, the replicas are independent and uniform. Expanding
$\xi_N$ in distinct Walsh characters and using their orthogonality
shows that
\begin{equation}
B_{1,N}(0)=V_N,
\qquad
B_{2,N}(0)=B_{3,N}(0)=0.
\label{eq:B-initial-values}
\end{equation}
Indeed, the first identity is \eqref{eq:V-N-def}. In each cross term,
at least one replica contains a nonconstant Walsh character appearing
only once, so its uniform expectation vanishes.

Apply Lemma~\ref{lem:stochastic-stability} to each of the three
observables in \eqref{eq:B1-def}--\eqref{eq:B3-def}. The observables
depend on at most four replicas. Every term produced by
\eqref{eq:stochastic-stability} is $N$ times the expectation of a
product of three covariance kernels. By H\"older's inequality,
replica symmetry, and Lemma~\ref{lem:cubic-kernel-moment},
\begin{equation}
\sup_{0\leq s\leq b^2}
|B_{j,N}'(s)|
\leq
C N\sup_{0\leq s\leq b^2}
\nu_s\left(|\xi_N(R_{1,2})|^3\right)
\leq
C N^{1-3p_\star/2},
\label{eq:B-derivative-bound}
\end{equation}
for $j=1,2,3$. Integrating from $0$ to $s$ and using
\eqref{eq:B-initial-values}, we obtain, uniformly for
$0\leq s\leq b^2$,
\begin{equation}
B_{1,N}(s)
=
V_N+O_{b,\xi}\bigl(N^{1-3p_\star/2}\bigr)
\label{eq:B1-asymptotic}
\end{equation}
and
\begin{equation}
B_{2,N}(s),\ B_{3,N}(s)
=
O_{b,\xi}\bigl(N^{1-3p_\star/2}\bigr).
\label{eq:B23-asymptotic}
\end{equation}
Substituting \eqref{eq:B1-asymptotic} and
\eqref{eq:B23-asymptotic} into \eqref{eq:A-N-derivative} gives
\begin{equation}
A_N'(s)
=
NV_N+O_{b,\xi}\bigl(N^{2-3p_\star/2}\bigr).
\label{eq:A-prime-asymptotic}
\end{equation}
Finally, $A_N(0)=0$ because every interaction character has positive
degree and hence zero mean under the uniform measure. Integrating
\eqref{eq:A-prime-asymptotic} from $0$ to $s$ proves
\eqref{eq:linear-response}.
\end{proof}

\subsection{Completion of the proof of Theorem~\ref{thm:quadratic-approximation}}
\label{subsec:complete-quadratic}

We now return to $D_N$ from \eqref{eq:D-N-beta0}.

First, Gaussian integration by parts gives
\begin{equation}
\frac{d}{ds}\E\log Z_N(\sqrt{s})
=
\frac N2\left(\xi_N(1)-A_N(s)\right).
\label{eq:mean-free-energy-derivative}
\end{equation}
On the other hand, from the definition of $J_N$,
\begin{equation}
\E\bigl[NJ_N(\sqrt{s})\bigr]
=
\frac{Ns}{2}\xi_N(1).
\label{eq:mean-J-s}
\end{equation}
Differentiating \eqref{eq:mean-J-s} with respect to $s$, subtracting
the resulting identity from \eqref{eq:mean-free-energy-derivative},
and using that both quantities vanish at $s=0$, we obtain
\begin{equation}
\E D_N
=
-\frac N2\int_0^{\beta_0^2}A_N(s)\,ds.
\label{eq:mean-D-exact}
\end{equation}
By Proposition~\ref{prop:linear-response} and
\eqref{eq:V-N-order},
\begin{equation}
|\E D_N|
\leq
C_{b,\xi}N^{2-p_\star}.
\label{eq:mean-D-bound}
\end{equation}

It remains to control the variance. For $p\in\cP$ and
$I\in\binom{[N]}p$, differentiation with respect to the corresponding
Gaussian coupling gives
\begin{equation}
\frac{\partial D_N}{\partial g_{p,I}}
=
\frac{\beta_0\gamma_p}{N^{(p-1)/2}}
\langle\sigma_I\rangle_{\beta_0^2}
-
\frac{\beta_0^2\gamma_p^2}{N^{p-1}}g_{p,I}.
\label{eq:D-gradient}
\end{equation}
Here $\langle\cdot\rangle_{\beta_0^2}$ is the same Gibbs expectation
as $\langle\cdot\rangle_{\beta_0}$ in the original inverse-temperature
notation. The Gaussian Poincar\'e inequality therefore gives
\begin{equation}
\var(D_N)
\leq
\sum_{p\in\cP}
\sum_{I\in\binom{[N]}p}
\E\left(
\frac{\partial D_N}{\partial g_{p,I}}
\right)^2.
\label{eq:Poincare-D}
\end{equation}
For each $(p,I)$, Gaussian integration by parts gives
\begin{equation}
\E\left[
g_{p,I}\langle\sigma_I\rangle_{\beta_0^2}
\right]
=
\frac{\beta_0\gamma_p}{N^{(p-1)/2}}
\E\left[
1-\langle\sigma_I\rangle_{\beta_0^2}^2
\right].
\label{eq:IBP-local-magnetization}
\end{equation}
Expanding the square in \eqref{eq:Poincare-D}, using
\eqref{eq:IBP-local-magnetization}, and recalling
\eqref{eq:A-N-positive} and \eqref{eq:V-N-def}, we obtain
\begin{align}
\var(D_N)
\leq{}&
\beta_0^2N A_N(\beta_0^2)
-
\beta_0^4N^2V_N
\notag\\
&+
2\beta_0^4
\sum_{p\in\cP}
\frac{\gamma_p^4}{N^{2p-2}}
\sum_{I\in\binom{[N]}p}
\E\langle\sigma_I\rangle_{\beta_0^2}^2.
\label{eq:variance-D-expanded}
\end{align}
Proposition~\ref{prop:linear-response} gives the cancellation
\begin{equation}
\beta_0^2N A_N(\beta_0^2)
-
\beta_0^4N^2V_N
=
O_{b,\xi}\bigl(N^{3-3p_\star/2}\bigr).
\label{eq:variance-leading-cancellation}
\end{equation}
For the last term in \eqref{eq:variance-D-expanded}, note that
\begin{align}
&
\sum_{p\in\cP}
\frac{\gamma_p^4}{N^{2p-2}}
\sum_{I\in\binom{[N]}p}
\E\langle\sigma_I\rangle_{\beta_0^2}^2
\notag\\
&\qquad\leq
\left(
\max_{p\in\cP}\frac{\gamma_p^2}{N^{p-1}}
\right)
\sum_{p\in\cP}
\frac{\gamma_p^2}{N^{p-1}}
\sum_{I\in\binom{[N]}p}
\E\langle\sigma_I\rangle_{\beta_0^2}^2
\notag\\
&\qquad=
\left(
\max_{p\in\cP}\frac{\gamma_p^2}{N^{p-1}}
\right)
N A_N(\beta_0^2).
\label{eq:last-Poincare-term}
\end{align}
By Proposition~\ref{prop:linear-response} and
\eqref{eq:V-N-order},
\[
A_N(\beta_0^2)=O_{b,\xi}(N^{1-p_\star}),
\]
while
\[
\max_{p\in\cP}\frac{\gamma_p^2}{N^{p-1}}
=O_\xi(N^{1-p_\star}).
\]
Hence the right-hand side of \eqref{eq:last-Poincare-term} is
$O_{b,\xi}(N^{3-2p_\star})$, which is smaller than
$N^{3-3p_\star/2}$ because $p_\star>2$. Combining this with
\eqref{eq:variance-leading-cancellation} yields
\begin{equation}
\var(D_N)
\leq
C_{b,\xi}N^{3-3p_\star/2}.
\label{eq:variance-D-final}
\end{equation}

Finally, by \eqref{eq:scaled-D-N},
\begin{align}
&\E\left|
N^{p_\star/2}
\bigl(F_N(\beta_0)-J_N(\beta_0)\bigr)
\right|^2
\notag\\
&\qquad=
N^{p_\star-2}
\left(
\var(D_N)+(\E D_N)^2
\right)
\notag\\
&\qquad\leq
C_{b,\xi}
\left(
N^{1-p_\star/2}+N^{2-p_\star}
\right)
\longrightarrow0,
\label{eq:quadratic-approximation-complete}
\end{align}
where we used~\eqref{eq:mean-D-bound} and~\eqref{eq:variance-D-final}. The constant depends only on $b$ and the fixed mixture. Since $p_\star\geq3$ and $N\geq1$, one has $N^{2-p_\star}\leq N^{1-p_\star/2}$. Taking the supremum over the arbitrary evaluation point $\beta_0\in[0,b]$ therefore proves~\eqref{eq:quadratic-approximation-rate}. At $\beta_0=0$ both $F_N$ and $J_N$ vanish identically. The limit~\eqref{eq:quadratic-approximation-limit} follows because $1-p_\star/2<0$. This completes the proof of Theorem~\ref{thm:quadratic-approximation}.







\section{Proof of the central limit theorem}
\label{sec:clt-proof}
The quadratic approximation reduces the free-energy fluctuation to an independent sum of centered Gaussian squares. Only the smallest active degree contributes at the stated scale.

\begin{proof}[Proof of Theorem~\ref{thm:main-clt}]
Fix $0<\beta<\beta_c(\xi)$ and choose $b$ with $\beta\leq b<\beta_c(\xi)$. By definition,
\begin{equation}
N^{p_\star/2}\bigl(J_N(\beta)-\E J_N(\beta)\bigr)
=\frac{\beta^2}{2}\sum_{p\in\cP}\gamma_p^2N^{p_\star/2-p}
\sum_{I\in\binom{[N]}p}(g_{p,I}^2-1)\,.
\label{eq:quadratic-clt-decomposition}
\end{equation}
For $p=p_\star$, write $m_N=\binom N{p_\star}$. The $m_N$ summands $g_{p_\star,I}^2-1$ are independent and identically distributed, with mean zero and variance two. Since $m_N\to\infty$, the classical independent-sum central limit theorem gives
\begin{equation}
\frac1{\sqrt{2m_N}}\sum_{I\in\binom{[N]}{p_\star}}(g_{p_\star,I}^2-1)
\Longrightarrow\mathcal N(0,1)\,.
\label{eq:independent-square-clt}
\end{equation}
Moreover, $N^{-p_\star/2}\sqrt{2m_N}\to\sqrt{2/p_\star!}$. Hence the contribution of $p_\star$ in~\eqref{eq:quadratic-clt-decomposition} converges to the centered Gaussian law with variance $\beta^4\gamma_{p_\star}^4/(2p_\star!)$.

For each $p>p_\star$, the variance of its contribution in~\eqref{eq:quadratic-clt-decomposition} equals
\begin{equation}
\frac{\beta^4\gamma_p^4}{2}\,N^{p_\star-2p}\binom Np
=O(N^{p_\star-p})\longrightarrow0\,.
\label{eq:higher-degree-vanishing}
\end{equation}
There are finitely many active degrees, so their total contribution tends to zero in $L^2$. Theorem~\ref{thm:quadratic-approximation} also gives $N^{p_\star/2}(F_N(\beta)-J_N(\beta))\to0$ in $L^2$. Combining these facts with the exact identity $\E J_N(\beta)=\beta^2\xi_N(1)/2$ and Slutsky's theorem proves~\eqref{eq:main-clt}.
\end{proof}

\begin{bibdiv}
\begin{biblist}
\bib{ALR1987}{article}{
  author={Aizenman, Michael},
  author={Lebowitz, Joel L.},
  author={Ruelle, David},
  title={Some rigorous results on the Sherrington--Kirkpatrick spin glass model},
  journal={Communications in Mathematical Physics},
  volume={112},
  date={1987},
  pages={3--20},
  note={\href{https://www.ihes.fr/~ruelle/PUBLICATIONS/[90].pdf}{Full text}},
}
\bib{BanerjeeBelius2021}{article}{
  author={Banerjee, Debapratim},
  author={Belius, David},
  title={Fluctuations of the free energy of the mixed $p$-spin mean field spin glass model},
  date={2021},
  note={\href{https://arxiv.org/abs/2108.03109}{arXiv:2108.03109}},
}
\bib{BardinaEtAl2004}{article}{
  author={Bardina, Xavier},
  author={M\'arquez-Carreras, David},
  author={Rovira, Carles},
  author={Tindel, Samy},
  title={The $p$-Spin Interaction Model with External Field},
  journal={Potential Analysis},
  volume={21},
  number={4},
  date={2004},
  pages={311--362},
  doi={10.1023/B:POTA.0000034325.04634.f5},
}
\bib{BovierKurkovaLowe2002}{article}{
  author={Bovier, Anton},
  author={Kurkova, Irina},
  author={L\"owe, Matthias},
  title={Fluctuations of the free energy in the REM and the $p$-spin SK models},
  journal={The Annals of Probability},
  volume={30},
  number={2},
  date={2002},
  pages={605--651},
  doi={10.1214/aop/1023481004},
}
\bib{BovierSchertzer2024}{article}{
  author={Bovier, Anton},
  author={Schertzer, Adrien},
  title={Fluctuations of the free energy in $p$-spin SK models on two scales},
  journal={Probability Theory and Related Fields},
  volume={189},
  date={2024},
  pages={771--810},
  doi={10.1007/s00440-024-01296-y},
}
\bib{Chatterjee2009}{article}{
  author={Chatterjee, Sourav},
  title={Disorder chaos and multiple valleys in spin glasses},
  date={2009},
  note={\href{https://arxiv.org/abs/0907.3381}{arXiv:0907.3381}},
}
\bib{ChenDeyPanchenko2017}{article}{
  author={Chen, Wei-Kuo},
  author={Dey, Partha},
  author={Panchenko, Dmitry},
  title={Fluctuations of the free energy in the mixed $p$-spin models with external field},
  journal={Probability Theory and Related Fields},
  volume={168},
  number={1--2},
  date={2017},
  pages={41--53},
  doi={10.1007/s00440-016-0705-5},
}
\bib{ChenHighTemperature}{article}{
 author={Chen, Wei-Kuo},
 title={Phase transition in the spiked random tensor with Rademacher prior},
 journal={The Annals of Statistics},
 volume={47}, number={5}, date={2019}, pages={2734--2756},
 doi={10.1214/18-AOS1763},
}
\bib{CometsNeveu1995}{article}{
  author={Comets, Francis},
  author={Neveu, Jacques},
  title={The Sherrington--Kirkpatrick model of spin glasses and stochastic calculus: The high temperature case},
  journal={Communications in Mathematical Physics},
  volume={166},
  number={3},
  date={1995},
  pages={549--564},
  doi={10.1007/BF02099887},
}
\bib{DeyKang2026}{article}{
  author={Dey, Partha S.},
  author={Kang, Taegu},
  title={Fluctuations for the Sherrington--Kirkpatrick spin glass model near the critical temperature},
  date={2026},
  note={\href{https://arxiv.org/abs/2603.05636}{arXiv:2603.05636}},
}
\bib{DeyWuHypergraphs}{article}{
  author={Dey, Partha S.},
  author={Wu, Qiang},
  title={Hypergraph Counting and Fluctuations of Mixed $p$-Spin Glass Models at High Temperature},
  journal={Communications in Mathematical Physics},
  volume={407},
  date={2026},
  note={Article 230},
  doi={10.1007/s00220-026-05748-5},
  url={https://arxiv.org/abs/2212.14571v2},
}
\bib{DeyWuMultiSpecies2021}{article}{
  author={Dey, Partha S.},
  author={Wu, Qiang},
  title={Fluctuation Results for Multi-species Sherrington--Kirkpatrick Model in the Replica Symmetric Regime},
  journal={Journal of Statistical Physics},
  volume={185},
  number={3},
  date={2021},
  note={Article 22},
  doi={10.1007/s10955-021-02835-w},
}
\bib{DeyWuWeakField2023}{article}{
  author={Dey, Partha S.},
  author={Wu, Qiang},
  title={Mean Field Spin Glass Models Under Weak External Field},
  journal={Communications in Mathematical Physics},
  volume={402},
  number={2},
  date={2023},
  pages={1205--1258},
  doi={10.1007/s00220-023-04742-5},
}
\bib{GuerraToninelli2002}{article}{
  author={Guerra, Francesco},
  author={Toninelli, Fabio L.},
  title={Central limit theorem for fluctuations in the high temperature region of the Sherrington--Kirkpatrick spin glass model},
  journal={Journal of Mathematical Physics},
  volume={43},
  number={12},
  date={2002},
  pages={6224--6237},
  doi={10.1063/1.1515109},
}
\bib{KnopfelLowe2005}{article}{
  author={Kn\"opfel, Holger},
  author={L\"owe, Matthias},
  title={Fluctuations in a $p$-spin interaction model},
  journal={Annales de l'Institut Henri Poincar\'e, Probabilit\'es et Statistiques},
  volume={41},
  number={4},
  date={2005},
  pages={807--815},
  doi={10.1016/j.anihpb.2004.05.006},
}
\bib{Panchenko2014}{article}{
  author={Panchenko, Dmitry},
  title={The Parisi formula for mixed $p$-spin models},
  journal={The Annals of Probability},
  volume={42},
  number={3},
  date={2014},
  pages={946--958},
  doi={10.1214/12-AOP800},
}
\bib{Talagrand2000}{article}{
  author={Talagrand, Michel},
  title={Rigorous low-temperature results for the mean field $p$-spins interaction model},
  journal={Probability Theory and Related Fields},
  volume={117},
  number={3},
  date={2000},
  pages={303--360},
  doi={10.1007/s004400050009},
}
\bib{Talagrand2006}{article}{
  author={Talagrand, Michel},
  title={The Parisi formula},
  journal={Annals of Mathematics},
  volume={163},
  number={1},
  date={2006},
  pages={221--263},
  doi={10.4007/annals.2006.163.221},
}
\bib{Tindel2005}{article}{
  author={Tindel, Samy},
  title={On the stochastic calculus method for spins systems},
  journal={The Annals of Probability},
  volume={33},
  number={2},
  date={2005},
  pages={561--581},
  doi={10.1214/009117904000000919},
}
\end{biblist}
\end{bibdiv}
 

\end{document}
